\chapter{Grotere getallen vermenigvuldigen}\label{ch:g4:mult}

In jaar~3 vermenigvuldigden we met \'{e}\'{e}n cijfer
(\cref{ch:g3:mult}); dit hoofdstuk vermenigvuldigt met tweecijferige
getallen --- en legt de mysterieuze ``verschoven regel'' van de
kolommethode uit, die niets anders is dan een splitsing in
tientallen en eenheden.

\section{Vermenigvuldigen met tientallen}

\begin{proposition}[Keer 10, 20, 300\,\dots]\label{prop:g4:mult:tens}
Vermenigvuldigen met $10$ voegt \'{e}\'{e}n nul toe, met $100$ twee
nullen. Om met $20$ te vermenigvuldigen, vermenigvuldig met $2$, dan
met $10$:
\[
36 \times 20 = 36 \times 2 \times 10 = 72 \times 10 = 720 .
\]
Zo ook $14 \times 300 = 14 \times 3 \times 100 = 4\,200$.
\end{proposition}
\admitted

\section{De kolommethode met twee cijfers}

\begin{example}[Waarom de verschoven regel?]\label{ex:g4:mult:why}
Om $23 \times 12$ uit te rekenen, splits $12$ in $10 + 2$:
\[
23 \times 12 = 23 \times 2 + 23 \times 10 = 46 + 230 = 276 .
\]
De kolommethode schrijft precies deze twee producten, de ene onder
de andere --- de tweede naar links verschoven omdat het een getal
van \emph{tientallen} is:
\[
\begin{array}{r}
2\,3 \\
\times\ 1\,2 \\
\hline
4\,6 \\
2\,3\,\cdot \\
\hline
2\,7\,6
\end{array}
\]
(de punt markeert de eenhedenplaats van de verschoven regel, vaak
blanco gelaten).
\end{example}

\begin{omfigure}
\begin{tikzpicture}[scale=0.42]
  \fill[omDef!18] (0,0) rectangle (20,10);
  \fill[omThm!18] (0,10) rectangle (20,12);
  \fill[omDef!34] (20,0) rectangle (23,10);
  \fill[omThm!34] (20,10) rectangle (23,12);
  \draw[thick] (0,0) rectangle (23,12);
  \draw[thick] (20,0) -- (20,12);
  \draw[thick] (0,10) -- (23,10);
  \node[font=\small] at (10,5) {$20 \times 10 = 200$};
  \node[font=\small] at (10,11) {$20 \times 2 = 40$};
  \node[font=\small, rotate=90] at (21.5,5) {$3 \times 10 = 30$};
  \node[font=\small] at (21.5,11) {$6$};
  \node[below, font=\small] at (10,0) {$20$};
  \node[below, font=\small] at (21.5,0) {$3$};
  \node[left, font=\small] at (0,5) {$10$};
  \node[left, font=\small] at (0,11) {$2$};
\end{tikzpicture}

{\small Het rechthoekplaatje van $23 \times 12$: vier makkelijke
producten, $200 + 40 + 30 + 6 = 276$. De kolommethode groepeert ze
in twee regels: $46$ (de ``$\times 2$''-rij) en $230$ (de
``$\times 10$''-rij).}
\end{omfigure}

\begin{method}[Kolomsgewijs vermenigvuldigen met een tweecijferig getal]\label{met:g4:mult:column}
Om $457 \times 36$ uit te rekenen:
\begin{enumerate}
  \item eerste regel: $457 \times 6 = 2\,742$;
  \item tweede regel: $457 \times 3$ \emph{tientallen} $= 1\,371$
        tientallen --- schrijf $1\,371$ \'{e}\'{e}n plaats naar links
        verschoven;
  \item tel de twee regels op: $2\,742 + 13\,710 = 16\,452$;
  \item schat om te controleren: $457 \times 36 \approx 500 \times 36
        = 18\,000$? Beter: $450 \times 36$ is ongeveer $16\,000$ ---
        aannemelijk.
\end{enumerate}
\[
\begin{array}{r}
4\,5\,7 \\
\times\ \ \,3\,6 \\
\hline
2\,7\,4\,2 \\
1\,3\,7\,1\,\cdot \\
\hline
1\,6\,4\,5\,2
\end{array}
\]
\end{method}

\begin{example}[Mentale trucs]\label{ex:g4:mult:tricks}
\begin{itemize}
  \item \emph{Vriendelijk splitsen}: $18 \times 11 = 18 \times 10
        + 18 = 198$.
  \item \emph{Bijna-honderd gebruiken}: $25 \times 99 = 25 \times 100 -
        25 = 2\,475$.
  \item \emph{Verdubbelen en halveren}: $16 \times 35 = 8 \times 70 =
        560$ (\'{e}\'{e}n factor halveren en de andere verdubbelen houdt
        het product).
\end{itemize}
\end{example}

\section{Vermenigvuldigen in sommen}

\begin{example}\label{ex:g4:mult:problem}
Een school bestelt $28$ dozen van $145$ vellen papier.
\begin{enumerate}
  \item Bewerking: $145 \times 28$.
  \item Schatting: $150 \times 30 = 4\,500$, een beetje te veel aan
        beide kanten.
  \item Kolommen: $145 \times 8 = 1\,160$; $145 \times 2$ tientallen
        $= 2\,900$; totaal $4\,060$.
  \item Zin: de school ontvangt $4\,060$ vellen.
\end{enumerate}
\end{example}

\section{Oefeningen}

\begin{exercise}[$\star$]\label{exo:g4:mult:1}
Reken uit het hoofd: $47 \times 10$; \quad $23 \times 100$; \quad
$36 \times 20$; \quad $15 \times 300$; \quad $42 \times 50$.
\end{exercise}

\begin{exercise}[$\star$]\label{exo:g4:mult:2}
Reken $34 \times 21$ uit door te splitsen zoals in
\cref{ex:g4:mult:why} ($34 \times 21 = 34 \times 20 + 34$),
controleer dan in kolommen.
\end{exercise}

\begin{exercise}[$\star$]\label{exo:g4:mult:3}
Reken kolomsgewijs uit: $63 \times 24$; \quad $85 \times 47$; \quad
$126 \times 53$.
\end{exercise}

\begin{exercise}[$\star$]\label{exo:g4:mult:4}
Reken kolomsgewijs uit, met eerst een schatting: $208 \times 34$;
\quad $517 \times 62$.
\end{exercise}

\begin{exercise}[$\star$]\label{exo:g4:mult:5}
Reken uit met een truc uit \cref{ex:g4:mult:tricks}:
$45 \times 11$; \quad $32 \times 99$; \quad $24 \times 45$
(verdubbel en half --- twee keer als je wilt).
\end{exercise}

\begin{exercise}[$\star$]\label{exo:g4:mult:6}
Een bioscoop heeft $26$ rijen van $18$ stoelen. Hoeveel stoelen in
totaal?
\end{exercise}

\begin{exercise}[$\star$]\label{exo:g4:mult:7}
Een vrachtwagen vervoert $32$ pallets van $48$ kratten. Hoeveel
kratten? Geef eerst een schatting, dan het exacte antwoord.
\end{exercise}

\begin{exercise}[$\star$]\label{exo:g4:mult:8}
Zoek de fout: Ana rekent $57 \times 23$ als $57 \times 3 = 171$ en
$57 \times 2 = 114$, en telt dan $171 + 114 = 285$. Het juiste
antwoord is $1\,311$. Wat is ze vergeten?
\end{exercise}

\begin{exercise}[$\star\star$]\label{exo:g4:mult:9}
Een tuinman plant $14$ rijen van $35$ tulpen en $12$ rijen van $28$
narcissen. Hoeveel bloemen in totaal? (Twee producten, dan een som.)
\end{exercise}

\begin{exercise}[$\star\star$]\label{exo:g4:mult:10}
Elke dag gebruikt een bakkerij $67$ kg meel. Ongeveer hoeveel meel
is dat in een maand van $30$ dagen --- schat met afgeronde getallen,
reken dan precies.
\end{exercise}

\begin{exercise}[$\star\star$]\label{exo:g4:mult:11}
Leg uit, met een rechthoekplaatje zoals in dit hoofdstuk, waarom
$15 \times 12$ berekend kan worden als $15 \times 10 + 15 \times 2$.
Bedenk daarna een andere splitsing van $15 \times 12$ die ook werkt.
\end{exercise}
